\scanpage{353}
% @ head of sheet
\mspage{77}
% @ top, continuation of b, lines 1–4
$L_1$-$\unsure{\kappa|\pi}$-adapté dans $\underline{\mathrm{Fl}}\,\mathrm{Cat}$ (cf ci-dessous), de limite inductive $f$, on a $f\in W$ \unsure{sss|ssi} la sous-\inserted{catégorie pl.} $I_W$ des $i\in\mathrm{Ob}\,I$ tels que $f_i\in \unsure{w|W}$ est cofinale dans $I$.

% @ upper third, lines 5–8
b$'$) (Quand on a choisi, pour \unsure{tte} $\deleted{\unsure{x}}f\in\deleted{\unsure{O}}\mathrm{Fl}(\mathrm{Cat})$, une représentation $L_1$-$\kappa$-adaptée de $f$ comme $\underrightarrow{\lim}$) : Si $f\in W$, alors la \deleted{\unsure{cat.}}\inserted{sous-cat. pl.} $I_W$ \unsure{définie} dans b) est cofinale dans $I$.

% @ centre, lines 9–10
Rappelons ici la notion de syst. inductif $L_1$-$\kappa$-adapté. Il s'agit des conditions :

% @ centre, conditions 1–3
1\textsuperscript{o}) $I$ est \unsure{grde|grande} devant $\kappa$, et stable par $\underrightarrow{\lim}$ filtrantes de cardinal $\leq\kappa$.

2\textsuperscript{o}) Le foncteur $i\mapsto x_i:I\longrightarrow\underline{\mathrm{Fl}}(\mathrm{Cat})$ commute aux $\underrightarrow{\lim}$ en question.

3\textsuperscript{o}) Les $x_i$ sont dans $\underline{\mathrm{Fl}}(\mathrm{Cat})_\kappa\overset{\mathrm{def}}{=}\underline{\mathrm{Fl}}_\kappa(\mathrm{Cat})$.

% @ lower third, lines 16–20
Notons que dans le cas d'un ``\unsure{paratopos|prétopos} \unsure{paradisien|paradislien|paradisiaque}'' tel que $\mathrm{Cat}$ ou $\underline{\mathrm{Fl}}\,\mathrm{Cat}$, il y a une façon canonique, $\pi$ étant donné, de représenter un objet $f$ comme $\underrightarrow{\lim}$ d'un syst. inductif $L_1$-$\kappa$-adapté.
% @ foot, last two lines
Ici, \deleted{$\unsure{f}$}
\[
f:X\longrightarrow Y
\]
étant donné, on considère l'ens $I$ des
% archive sheet number: 352
